\chapter{Практика}

\section{Задание 1}

Сформируем сигнал с двоичной AM с амплитудами $S_1$ = 0.4 для единицы и $S_2$ = 0 для нуля.
Сделаем 2 отсчета сигнала в момент времени $t_1$ и $t_2$ со значениями $z(t_1) = 0.3, z(t_2) = 0.2$. \\

Для принятия решения о переданном бите необходимо вычислить условную плотность распределения сигнала при условии передачи $S_1$ и $S_2$ в
точках z. \\

Сравнив значения, полученные выше, мы получим значение переданного бита:

\begin{figure}[H]
    \centering
    \includegraphics{1.png}
    \caption{Результат работы РУ}
\end{figure}


\section{Задание 2}
Сделаем то же самое, но уже сами вычислим и добавим к сигналу шум. Для вычисления шума даны: $R_b = 100\frac{symbols}{s}$, $\frac{N_0}{2}=0.25*10^{-2}$.
Шум равен произведению спектральной плотности мощности и ширины полосы (битовой скорости). \\

Также проверим работу РУ при разных амплитудах символа, передающего единицу:

\begin{figure}[H]
    \centering
    \includegraphics{2.png}
    \caption{Результат работы РУ при разных амплитудах}
\end{figure}

На передатчике передавалась единица, но приемник смог дать правильный ответ только при повышенной амплитуде. Чем амплитуда выше относительно шума, тем выше точность работы РУ.

\section{Задание 3}

Примем BSPK символ и подадим на РУ. При наличии только двух точек в созвездии модуляции можно упростить логику РУ до 

$$z > \frac{s_1+s_2}{2} \rightarrow 1$$

Здесь z - отсчет принятого сигнала.

\begin{figure}[H]
    \centering
    \includegraphics{3.png}
    \caption{Результат работы РУ}
\end{figure}

\section{Задание 4}

Повторим действия из прошлого шага, но теперь сгенерируем 1000 символов и вычислим Bit Error Rate (BER) как отношение кол-ва ошибок к общему кол-ву переданных символов.

\begin{figure}[H]
    \centering
    \includegraphics{4.png}
    \caption{BER на приеме}
\end{figure}

Такая метрика говорит о том, что наше РУ дает 205 ошибок на 1000 отсчетов.

\section{Задание 5}

Проделаем шаги из прошлого шага, но теперь посмотрим на зависимость BER от амплитуды символа 
\begin{figure}[H]
    \centering
    \includegraphics{5.png}
    \caption{Зависимость BER от амплитуды}
\end{figure}

Кол-во ошибок уменьшается с ростом амплитуды, потому что шум становится малым относительно амплитуды символа и символ можно легче "узнать".

\endinput 